\section{引言}

电力、石化、冶金、交通、新能源等工业领域的重大基础设施通常长期运行于高温、高压、强电磁、粉尘、烟雾以及高空、狭窄等复杂环境，其设备与结构在长期服役过程中不可避免地出现老化、损伤与故障。及时、准确地开展巡检与状态评估，是保障工业设施安全运行和实现智能运维的重要基础。传统人工巡检虽然具有较强的现场判断能力，但在高危、难抵达和大范围场景下普遍存在劳动强度大、作业风险高、巡检效率有限以及结果受人员经验影响明显等问题。无人机（Unmanned Aerial Vehicle, UAV）具有机动灵活、非接触作业、视角可控以及可搭载多种传感器等优势，已经逐步应用于输电线路、桥梁、风电、光伏和厂区设备巡检\cite{shakhatreh2019survey,floreano2015science,kumar2012opportunities,cai2014survey,kendoul2012survey}。已有研究从无人机系统、计算机视觉、摄影测量与遥感等不同角度总结了相关技术的发展，并表明无人机已经成为工业设施自动化巡检的重要技术载体\cite{kanellakis2017survey,colomina2014unmanned,nex2014uav}。

随着无人机平台、机载传感器和人工智能技术的发展，工业巡检正在由人工操控和预设航线采集，逐步向自主感知、自动识别和自主决策方向演进。在电力和基础设施巡检领域，已有研究利用无人机图像开展线路部件和缺陷检测\cite{nguyen2018automatic,miao2020uav,zhouwq2024review}；在桥梁及其他基础设施检测中，无人机视觉系统已能够完成裂缝、剥落和锈蚀等表观病害识别\cite{morgenthal2014quality,ham2016visual,dorafshan2018bridge,spencer2019advances}；在工业设施内部，也已有无人机开展受限空间检测和自主导航的研究\cite{nikolic2013uav,ozaslan2017autonomous,petrlik2020robust}。这些研究证明了无人机巡检技术的可行性，但复杂工业场景与普通室外环境相比具有更加显著的环境和任务约束。

复杂工业场景中的厂房、管廊、储罐、塔架和大型设备通常具有钢结构密集、空间遮挡严重、纹理重复以及GNSS信号受遮挡和多径效应影响等特点，同时还可能存在强电磁干扰、光照剧烈变化、粉尘烟雾、雨雾以及动态人员和车辆等因素\cite{scaramuzza2014vision,stoecker2017review,lvyang2019sense}。从巡检对象看，裂纹、腐蚀、螺栓、销钉和局部过热等缺陷通常尺度较小，且真实工业缺陷样本具有明显的稀缺性、类别不平衡和长尾分布特征\cite{taox2021survey,nguyen2018automatic}。从任务层面看，工业设备通常分布范围广、巡检目标多，单架无人机的续航、算力和通信资源有限，多无人机协同还需要同时考虑任务分配、碰撞规避、通信和能量约束\cite{zeng2016wireless,mozaffari2019tutorial,zhang2019energy}。因此，真正意义上的复杂工业场景无人机自主巡检，并非简单地按照预设航线采集图像，而是需要无人机基于环境信息和设备状态，自主完成定位建图、巡检对象观测、异常识别、任务规划以及任务调整等连续过程。

围绕上述问题，本文将复杂工业场景无人机自主巡检划分为三个相互衔接的层次：首先是面向飞行安全和环境理解的自主感知与定位，其核心是解决复杂退化环境下“无人机在哪里、周围有什么”的问题；其次是面向巡检对象的缺陷检测与异常认知，其核心是解决“设备是否异常、异常发生在哪里以及异常程度如何”的问题；最后是面向巡检任务的自主规划与多无人机协同，其核心是解决“下一步应该如何巡检以及多架无人机如何协同完成任务”的问题。本文在梳理国内外相关研究的基础上，重点分析不同技术路线的适用条件与局限性，并进一步总结当前复杂工业场景无人机自主巡检中仍存在的关键问题和未来发展趋势。

\section{复杂工业场景下无人机自主感知与信息融合}

\subsection{复杂工业场景的感知需求与自主定位问题}

复杂工业场景的巡检对象既包括输电线路、桥梁、风机和光伏组件等大范围设施，也包括厂房内部管道、阀门、储罐、塔架和设备连接部件等局部目标，因此无人机既需要完成大范围场景的快速搜索，也需要在接近目标后进行近距离精细观测。现有研究表明，工业巡检中的感知任务大体可以归纳为三类：一是针对无人机自身飞行安全的环境感知，包括障碍物、地形以及动态目标的识别，为自主导航和避障提供信息；二是针对巡检对象的目标感知，包括设备部件和缺陷区域的检测、识别与定位；三是面向场景理解的几何与语义建图，为任务规划、变化检测和数字化运维提供空间基础\cite{kanellakis2017survey,yeum2015vision,cha2017deep,nex2014uav,tao2019digital}。在室内、地下或大型工业设施内部，由于GNSS信号弱甚至完全不可用，无人机还必须依赖机载传感器完成自身状态估计和环境建图\cite{ozaslan2017autonomous,petrlik2020robust,nil2012review}。

从具体传感器看，可见光相机具有成本低、信息丰富等特点，是工业缺陷识别最常用的数据来源；红外热像仪能够反映设备温度分布，对电气设备过热、光伏热斑等异常具有较好的识别能力\cite{jadin2012recent,usamentiaga2014infrared}；激光雷达能够提供高精度三维几何信息，对光照变化不敏感，适用于结构测量、定位建图和障碍物感知\cite{wujq2021review,zhang2014loam}；IMU与GNSS则为无人机运动状态估计提供基础观测量。不同传感器具有明显的互补性，单一传感器往往难以同时满足工业环境下的定位精度、实时性和鲁棒性要求，因此多传感器融合成为自主巡检感知系统的重要发展方向\cite{panq2003basic,wangj2004survey,hey2000multi}。

无人机自主定位通常以同步定位与建图（Simultaneous Localization and Mapping, SLAM）为核心。视觉SLAM以图像为主要输入，具有传感器成本低和信息丰富等优势，ORB-SLAM系列代表了特征法视觉SLAM的发展路线\cite{durrant2006simultaneous,cadena2016past,mur2015orb,mur2017orb,campos2021orb}；LSD-SLAM和SVO等方法则分别代表直接法和半直接法，在大尺度建图和高帧率视觉里程计方面具有一定优势。激光SLAM能够直接利用环境几何信息，LOAM、FAST-LIO2等方法通过激光与惯性信息融合提高了定位精度与实时性\cite{zhang2014loam,xu2022fast}；VINS-Mono等视觉惯性方法则利用相机和IMU之间的互补关系实现自主状态估计\cite{qin2018vins}。近年来，多传感器融合逐步由单纯的松耦合组合向紧耦合状态估计发展，直接利用原始点云或图像测量信息能够减少中间特征提取环节，并提升复杂环境下的状态估计效率\cite{gupta2022slam,bai2022fasterlio}。

对于工业巡检而言，不同定位方法并不存在绝对优劣，而是与具体作业环境和机载资源密切相关。视觉方法在光照稳定、纹理丰富的环境中具有较好的成本和计算优势，但容易受到弱纹理、强反光、动态遮挡和光照变化影响；激光方法具有较好的几何鲁棒性，但硬件成本和计算负荷相对较高；视觉--惯性及激光--惯性紧耦合方法能够进一步提升定位稳定性，但对传感器外参标定、时间同步以及机载算力提出了更高要求。针对这些问题，已有研究在GPS拒止和受限空间环境中开展了自主飞行验证，例如Ozaslan等将自主导航和三维建图技术用于水电压力钢管与隧道检测\cite{ozaslan2017autonomous}，Petrlík等则在地下受限环境中验证了鲁棒无人机自主系统\cite{petrlik2020robust}。这些工作说明复杂环境中的自主感知已经具备较好的技术基础，但距离工业现场长期、稳定和可靠运行仍存在一定差距。

总体来看，现有研究已经形成了视觉、激光、视觉--惯性和激光--惯性等多种自主感知路线，多模态融合显著提升了无人机在复杂环境中的定位与建图能力。然而，工业现场还存在钢结构密集、GNSS多径、强反光、粉尘烟雾、动态遮挡和强电磁干扰等综合退化因素，容易造成感知信息质量下降和状态估计漂移\cite{lvyang2019sense,stoecker2017review,petrlik2020robust}。因此，复杂工业场景下的自主感知不仅需要提高单一算法的精度，更需要具备对传感器退化和异常观测的识别能力，并能够在不同感知条件下进行自适应融合。

\subsection{多模态感知融合与SLAM研究}

多模态信息融合通常可以划分为数据级、特征级和决策级三种方式。数据级融合能够保留较为完整的原始信息，但对传感器时间同步和空间标定提出了较高要求；特征级融合能够在信息量和计算开销之间取得平衡，是视觉、激光和惯性组合的重要实现形式；决策级融合则具有较强的模块独立性，适用于异构传感器之间的结果互证和冗余备份\cite{panq2003basic,hey2000multi}。在自主导航领域，多传感器融合的目标已经由简单地叠加多个传感器观测，逐渐发展为对状态估计、地图构建以及环境语义进行联合建模。Gupta和Fernando\cite{gupta2022slam}对无人机SLAM和多传感器融合研究进行了系统总结，表明视觉、激光、惯性等多源信息的深度融合是提升无人机自主感知鲁棒性的主要方向之一。

近年来，激光--惯性融合方法受到广泛关注。FAST-LIO2\cite{xu2022fast}直接利用原始点云构建残差，并结合IMU预测与动态地图结构实现高频率状态估计；Faster-LIO\cite{bai2022fasterlio}进一步从数据结构和并行计算角度降低了实时建图的计算开销。这类紧耦合方法避免了复杂的人工特征提取过程，在计算资源有限的机载平台上具有明显优势。与此同时，视觉--惯性方法在轻量化无人机平台中仍具有重要价值，Loianno等利用单目相机与IMU实现了高速飞行\cite{loianno2017estimation}，Faessler等验证了视觉自主飞行与三维建图的可行性\cite{faessler2016autonomous}。对于工业巡检无人机而言，视觉和激光并不是互相替代的关系，而是可以根据具体任务形成互补：视觉主要提供丰富的纹理和语义信息，激光主要提供稳定的空间几何信息，IMU则承担高频运动状态约束。

随着深度学习的发展，感知融合进一步从单纯的几何状态估计扩展至语义理解。视觉、激光和惯性信息能够共同支撑三维场景重建，而语义信息则可以进一步赋予地图中的设备、部件和缺陷区域明确的类别属性，为后续巡检目标定位和任务规划提供依据\cite{huangzx2023survey,tao2019digital}。主动SLAM与下一最佳视点规划等研究则开始考虑“为了获得更有价值的观测，应当如何运动”的问题\cite{placed2023active,bircher2016receding}。这说明自主巡检感知正在从被动获取环境信息向任务驱动的主动感知方向发展。

总体而言，多模态融合已经从“提高定位精度”的单一目标逐步拓展到“同时服务环境理解、目标认知和任务规划”的系统层面。但现有研究仍较多集中于算法本身的定位精度和计算效率，对于工业现场特殊退化因素的长期影响、不同传感器失效情况下的在线重构以及感知结果对巡检任务决策的直接反馈研究不足。因此，未来复杂工业场景无人机感知需要进一步加强退化感知、跨模态互补以及任务驱动的主动感知，使感知系统真正成为自主巡检闭环的一部分。

\section{工业设备缺陷检测与异常认知}

\subsection{基于视觉与多模态信息的工业缺陷检测}

工业设备异常具有明显的多样性，既包括裂纹、锈蚀、剥落、断股等结构表观缺陷，也包括部件缺失、松动、破损以及异物悬挂等状态异常，还包括电气设备过热、光伏热斑等热学异常以及烟火、树障和人员车辆入侵等环境异常\cite{taox2021survey,nguyen2018automatic,miao2020uav}。从认知粒度看，工业巡检通常包括图像级异常分类、目标级缺陷检测和像素级缺陷分割三个层次\cite{taox2021survey}。其中，目标检测主要解决“异常在哪里”的问题，像素级分割则进一步刻画缺陷的具体形态。MVTec AD等公开数据集推动了工业异常检测从单一案例验证走向标准化评测\cite{bergmann2019mvtec}。

在监督式视觉检测方面，卷积神经网络长期占据重要地位。两阶段检测方法以Faster R-CNN\cite{ren2015faster}等为代表，通常具有较好的检测精度；单阶段检测方法以YOLO系列、SSD和RetinaNet等为代表，在实时性和计算效率方面更具优势\cite{redmon2016you,redmon2017yolo9000,redmon2018yolov3,bochkovskiy2020yolov4,wang2023yolov7,liu2016ssd,lin2017focal}。近年来，DETR、Deformable DETR以及视觉Transformer等方法进一步增强了复杂场景中的特征建模能力\cite{carion2020end,zhu2021deformable,dosovitskiy2021image,liu2021swin,vaswani2017attention}。在机载边缘计算条件下，MobileNet和EfficientNet等轻量化网络能够在精度与计算开销之间取得一定平衡\cite{howard2017mobilenets,tan2019efficientnet}。

上述方法已经在桥梁、输电线路和工业设备中得到广泛应用。例如，Cha等利用深度网络实现混凝土裂缝检测\cite{cha2017deep}，Yeum和Dyke开展了桥梁裂缝视觉自动检测\cite{yeum2015vision}，Liu等将无人机图像与三维重建结合用于桥墩裂缝评估\cite{liu2020image}，Tao等针对航拍场景中的绝缘子缺陷开展了目标检测研究\cite{tao2020detection}。这些研究说明监督式深度学习已经能够较好解决“已知设备+已知缺陷类型”的检测问题，并且可以通过迁移学习、轻量化和多尺度特征建模提升工程适用性。

但是，复杂工业巡检与标准图像检测任务之间仍存在明显差异。首先，无人机与巡检对象之间的距离和视角不断变化，使缺陷在图像中的尺度变化更加明显；其次，裂纹、腐蚀和螺栓等目标通常尺寸较小，且容易受到背景纹理、反光和运动模糊影响；再次，工业设备的类别多、型号复杂，真实缺陷样本往往数量有限，导致模型容易依赖特定设备和特定场景的数据分布\cite{taox2021survey,zhouwq2024review}。因此，单纯追求公开数据集上的检测精度并不能完全代表复杂工业自主巡检的实际能力，检测模型的跨设备、跨场景泛化能力以及在机载平台上的实时运行能力同样值得重视。

对于具有明显热异常的设备，红外热成像提供了区别于可见光的重要信息。已有研究已经将红外技术用于电气设备故障诊断和无损检测\cite{jadin2012recent,usamentiaga2014infrared}。可见光图像能够反映裂纹、锈蚀和表面破损等结构信息，而红外图像能够反映温度场变化，两者具有明显的互补性。因此，将多模态数据进行联合分析，有望降低单一视觉模态在光照变化和低对比度条件下的检测风险，并为后续异常等级评估和智能预警提供更加丰富的信息基础。

\subsection{小样本、无监督异常检测与基础模型认知}

工业缺陷检测面临的核心困难之一是缺陷样本稀缺。与通用视觉任务不同，工业生产通常以正常状态为主，真实缺陷出现频率低，而且不同设备和不同工况之间缺陷类型存在明显差异，因此大量获取高质量缺陷标注具有较高成本。针对这一问题，研究逐渐由依赖人工标注的监督学习转向无监督和自监督异常检测\cite{taox2021survey,pang2021deep}。

早期无监督异常检测主要依赖重建和分布建模思想，例如VAE和GAN为基于生成模型的异常检测提供了基础\cite{kingma2014auto,goodfellow2014generative}，f-AnoGAN等方法进一步利用生成模型隐空间进行异常定位；Deep SVDD则通过学习正常样本的特征边界实现单类异常检测\cite{ruff2018deep,fangan2019}。随着预训练视觉特征的发展，基于特征分布与记忆库的方法逐渐成为重要路线。SPADE利用深层特征对应关系进行异常检测，PaDiM通过高斯分布描述图像块特征，PatchCore\cite{roth2022towards}则利用核心集采样和特征记忆库实现高效异常定位；EfficientAD\cite{batzner2024efficientad}进一步强调在线检测的计算效率。与此同时，CutPaste和DRAEM等自监督方法通过构造人工异常训练模型，从而减少对真实缺陷标注的依赖\cite{zavrtanik2021draem}。这些方法对于工业巡检的意义在于，它们能够在缺陷数据有限的情况下利用大量正常样本学习设备的正常状态。

然而，无监督异常检测也并非能够直接解决真实工业场景中的全部问题。其一，许多方法仍建立在训练数据与测试数据具有较接近设备形态和环境分布的假设上，当设备型号、拍摄距离和光照条件发生较大变化时，正常样本本身的分布也会发生变化；其二，工业现场存在的异常并不一定表现为明显的局部纹理差异，一些由温度变化、结构变形或设备组合状态引起的异常需要结合多模态和上下文信息进行判断；其三，实际巡检不仅需要检测异常，还需要进一步确定异常等级、空间位置、历史变化趋势和处置优先级。因此，未来的工业异常认知需要从单张图像异常检测进一步发展到多模态、时序化和知识辅助的状态认知。

视觉基础模型的发展为这一问题提供了新的研究路径。CLIP通过大规模图像--文本对比学习获得了较强的跨场景视觉表征能力\cite{radford2021learning}，WinCLIP进一步探索了零样本和少样本工业异常检测\cite{jeong2023winclip}；SAM提供了通用的可提示分割能力\cite{kirillov2023segment}，Grounding DINO支持开放词汇目标检测，DINOv2则能够提供较为鲁棒的无标注视觉特征\cite{liu2024grounding,oquab2024dinov2}。这些基础模型使异常认知开始由固定类别检测向开放词汇和少样本认知转变，为工业场景中的长尾缺陷和未知异常提供了新的技术基础。

在智能预警层面，异常检测结果还需要与设备台账、历史巡检记录以及设备运行参数进行关联，才能从单次“发现异常”进一步发展为状态评估和趋势预测。将无人机巡检数据与数字孪生模型结合，可以在统一空间中关联设备结构、历史状态与异常信息，为缺陷演化分析和预测性维护提供基础\cite{tao2019digital}。因此，工业异常认知的研究正在形成从“缺陷检测”向“状态理解”的演进路线，即由已知缺陷的监督识别逐步扩展到小样本、无监督和开放世界异常认知，并进一步与多模态信息和数字孪生结合，形成“检测--认知--预警”的闭环。

\section{无人机自主巡检规划与多无人机协同}

\subsection{自主巡检任务规划与动态路径重规划}

自主巡检任务规划的核心是在满足飞行安全、目标覆盖、观测质量、续航以及通信等约束条件下，确定无人机的飞行路径和作业顺序。传统无人机路径规划研究可以分为图搜索、随机采样、智能优化以及强化学习等路线\cite{wangq2019overview,liuqj2023survey}。其中，Dijkstra和A*为图搜索方法提供了经典基础\cite{dijkstra1959note,hart1968formal}；PRM、RRT以及RRT*适用于复杂高维空间中的可行路径搜索\cite{kavraki1996probabilistic,lavalle2001randomized,karaman2011sampling}；人工势场、动态窗口和互惠速度障碍等方法则主要用于局部在线避障\cite{khatib1986real,fox1997dynamic,van2011reciprocal}。这些方法在一般机器人导航中已经较为成熟，但工业巡检与普通导航存在明显区别：无人机不是单纯从起点到达终点，而是需要对设备表面或目标区域进行完整观测，因此覆盖率、观测距离和视角质量往往与路径长度具有同等重要性。

针对大范围巡检任务，覆盖路径规划（Coverage Path Planning, CPP）成为重要研究方向。Choset以及Galceran等对覆盖规划的基本理论和算法进行了系统总结\cite{choset2001coverage,galceran2013survey}，相关研究也进一步面向无人机进行了扩展。与普通最短路径问题相比，巡检覆盖规划需要综合考虑设备几何结构、障碍物布局、传感器视场以及无人机运动能力，因此越来越强调几何路径与观测任务之间的耦合。另一方面，当环境信息不断变化或存在动态目标时，单纯的离线全局路径难以持续适应。D*和D* Lite等增量式搜索方法能够根据环境更新快速进行路径修正\cite{stentz1994optimal,koenig2002d}，EGO-Planner等局部轨迹优化方法进一步降低了在线规划对复杂地图结构的依赖\cite{zhou2022ego}；Bircher等提出的下一最佳视点规划则从信息获取角度考虑自主探索和观测价值\cite{bircher2016receding}。

为了满足无人机动力学可行性，规划结果还需要与轨迹优化和飞行控制相结合。Mellinger和Kumar提出的最小snap轨迹生成方法能够生成适用于四旋翼飞行的动力学可行轨迹\cite{mellinger2011minimum}，Usenko等利用B样条实现了实时轨迹重规划\cite{usenko2017real}。这些研究说明，自主巡检规划正在由“几何路径搜索”逐步向“任务约束下的动态轨迹生成”转变。然而，对于工业巡检而言，现有方法仍较少同时考虑设备状态、异常程度、观测质量和剩余能量等任务信息。实际系统中，巡检任务往往不是一次性规划完成，而应根据当前的检测结果决定是否靠近目标复查、是否改变巡检顺序以及是否触发临时任务。因此，面向工业设备状态的闭环自主规划仍是值得深入研究的问题。

总体而言，现有规划研究已经建立了较为完善的从全局搜索、覆盖规划到局部重规划和动力学轨迹生成的技术体系，但许多研究主要以路径长度、时间或避障安全为目标，而真实工业巡检还要求规划结果服务于“发现异常--重新观测--确认状态”的任务闭环。未来需要进一步推动感知结果与任务规划之间的信息反馈，使无人机能够根据设备异常等级和环境变化动态调整巡检路径，而不是长期依赖预设航线。

\subsection{多无人机任务分配与协同决策}

面对大范围、多目标的工业设施巡检任务，单架无人机容易受到续航、载荷和作业时间等因素限制，多无人机协同能够通过任务并行提高巡检效率，并在单机故障时提升系统的任务鲁棒性。多机器人任务分配研究已经形成较成熟的理论框架\cite{gerkey2004formal}，其中市场机制和分布式拍卖是常见的任务分配方法，CBBA通过一致性协议实现分布式任务分配，在任务收益和通信开销之间取得了较好的平衡\cite{choi2009consensus}。国内也已有研究针对无人机集群任务规划和任务分配进行了系统总结，并提出粒子群、博弈论等方法解决多目标和异构任务分配问题\cite{jiagw2021review,duyh2020survey,jiayn2020recent,zhaiz2023market,zhangrp2022hybrid,jukai2022task}。

随着无人机数量增加，任务分配逐渐由集中式决策向分布式协同演进。一致性控制理论为多无人机协同提供了基本数学框架\cite{olfati2004consensus,ren2005consensus,olfati2006flocking}，而多智能体强化学习则进一步提供了基于数据的协同决策机制。MADDPG、QMIX和MAPPO等方法能够在多智能体环境中学习协作策略\cite{lowe2017multi,rashid2018qmix,yu2022surprising}。近年来，去中心化规划和在线协同探索开始在真实飞行平台上得到验证。Zhou等实现了仅依靠机载感知和局部通信的多无人机去中心化自主飞行，并进一步提出RACER等系统开展快速协同探索\cite{zhou2021decentralized,zhou2022swarm,zhou2023racer}。这些研究说明多无人机协同已经逐步从仿真环境向真实飞行实验发展。

然而，多无人机协同巡检面临的实际约束比一般集群飞行更加复杂。首先，工业厂区的无线通信容易受到建筑物、钢结构和电磁环境影响，通信链路可能产生带宽限制、延迟甚至中断\cite{zeng2016wireless,mozaffari2019tutorial}；其次，不同无人机可能搭载不同传感器和载荷，异构平台之间的任务收益和续航能力存在差异；再次，巡检任务往往会随着设备异常状态不断变化，任务分配必须具备在线重分配能力；此外，多机协同还需要同时考虑能量和计算资源约束\cite{zhang2019energy,zhou2018computation}。因此，现有“理想通信+静态任务”的多机协同研究还不能完全覆盖真实工业巡检需求。

总体来看，多无人机巡检正在从集中式离线任务分配向分布式在线协同决策发展，但真实工业环境中的通信受限、异构平台、动态任务以及安全互避等问题仍有待系统解决。未来，多无人机协同需要与环境感知和异常认知进一步耦合，使任务分配不再只根据空间距离和剩余能量进行，而能够根据设备风险等级、异常优先级和任务紧急程度动态决定无人机的协同方式，从而形成由巡检状态驱动的自主任务决策机制。

\section{总结与展望}

\subsection{现有研究存在的问题}

综合上述研究可以看出，复杂工业场景无人机自主巡检已经在环境感知、缺陷检测、路径规划以及多无人机协同等方面形成较为完整的技术体系，但距离复杂工业环境下长期、稳定、闭环运行的自主巡检仍存在以下几个方面的问题。

（1）复杂退化环境下自主感知的鲁棒性仍然不足。现有视觉、激光以及多模态SLAM方法已经能够在部分GNSS拒止和受限空间中实现自主定位，但对于强反光、弱纹理、粉尘烟雾、动态遮挡和复杂电磁干扰等综合退化条件，传感器观测质量可能显著下降，现有方法对传感器异常和失效的在线识别及自适应融合能力仍有限\cite{cadena2016past,lvyang2019sense,wujq2021review,petrlik2020robust}。

（2）工业异常认知仍存在跨场景泛化和未知异常识别能力不足的问题。监督式检测方法在特定设备和特定缺陷类别上能够取得较高精度，但对标注数据依赖较强；无监督异常检测能够缓解缺陷标注需求，却仍容易受到设备型号、拍摄环境和正常状态分布变化的影响\cite{taox2021survey,roth2022towards,batzner2024efficientad}。基础模型为少样本、零样本和开放词汇认知提供了新的可能，但其在真实工业场景中的可靠性、领域适应性以及异常等级判断能力仍缺乏充分验证\cite{jeong2023winclip,kirillov2023segment,liu2024grounding,oquab2024dinov2}。

（3）现有自主巡检规划与设备状态之间耦合不足。传统路径规划主要关注路径长度、避障和运动学可行性，而实际工业巡检还需要考虑巡检对象的重要程度、异常等级、观测质量、剩余能量和通信状态。已有研究虽然开始探索主动感知和下一最佳视点规划，但感知结果实时驱动任务重规划的研究仍不充分\cite{bircher2016receding,placed2023active}。

（4）多无人机协同在真实工业环境中的实用化程度仍有限。现有多机任务分配和协同决策研究已经从集中式方法发展到分布式和学习型方法，但在通信受限、异构平台、动态任务以及安全互避等约束下的长期稳定运行仍缺少大规模真实工业场景验证\cite{choi2009consensus,biwh2024review,xiangjw2022key,zhou2023racer}。

（5）感知、认知与决策之间尚未形成真正闭环。当前大量研究仍围绕单一模块展开，例如SLAM强调定位精度，缺陷检测强调识别精度，路径规划强调轨迹质量，多无人机协同强调任务分配效率，不同模块之间的信息传递和任务反馈仍相对有限。对于真正的自主巡检系统，无人机应能够根据“环境状态—设备状态—任务状态”的变化持续调整自身行为，因此需要进一步加强感知、异常认知和任务决策的系统级融合。

\subsection{未来发展趋势}

针对上述问题，未来复杂工业场景无人机自主巡检可以从以下几个方向进一步发展。

（1）面向退化环境的多模态鲁棒感知。未来需要从单纯提升定位精度转向提高感知系统对传感器退化和异常观测的容忍能力，通过视觉、激光、惯性、红外等多源信息的互补以及在线可信度评估，实现不同感知模式之间的自适应切换与融合，从而提高无人机在GNSS拒止、强反光、粉尘烟雾和动态遮挡环境中的长期自主运行能力。

（2）基础模型驱动的开放世界异常认知。CLIP、SAM、Grounding DINO和DINOv2等基础模型为少样本、开放词汇和未知异常识别提供了新的技术基础\cite{radford2021learning,kirillov2023segment,liu2024grounding,oquab2024dinov2}。未来可以进一步结合工业领域知识、设备结构信息和历史巡检数据，使无人机能够从“识别预先定义的缺陷”逐渐转向“发现和解释未知异常”，并进一步实现异常等级评估和智能预警。

（3）感知驱动的主动巡检与状态闭环。未来巡检规划将由固定航线逐渐转向根据环境和设备状态动态生成任务。无人机可以利用异常检测结果判断是否需要靠近目标、改变观测视角、增加复检次数或临时改变巡检顺序，并通过主动感知和下一最佳视点规划提高有限飞行时间内的信息获取效率\cite{bircher2016receding,placed2023active}。进一步结合设备数字孪生，可以将单次巡检结果与历史状态关联，支持缺陷演化分析和预测性维护\cite{tao2019digital}。

（4）通信受限条件下的多无人机自主协同。未来多无人机巡检需要从静态任务分配发展到动态、分布式和风险驱动的协同决策，将设备异常等级、任务紧急程度、剩余能量和通信状态共同纳入任务分配过程，并结合多智能体强化学习、分布式优化和端边云协同计算，提高大范围工业设施巡检的效率与系统鲁棒性\cite{yu2022surprising,zhou2023racer,zhang2019energy,zhou2018computation}。

总体而言，复杂工业场景无人机自主巡检的发展方向并不是单纯追求某一种感知、检测或规划算法的性能提升，而是逐步形成“多模态感知—异常认知—自主决策—任务执行—状态反馈”的闭环智能系统。随着基础模型、多模态感知、主动规划、多无人机协同和数字孪生等技术进一步融合，无人机巡检有望由当前的“自动化数据采集”逐步发展为面向工业设备状态的“自主感知、主动巡检与智能运维”系统。
